\bookchapter{Chapter 31: Getting Started with Short Video Shooting and Content Creation}

\chapterintro{In 2026, short videos are no longer “toys for young people” — they've become the fundamental way everyone gets information and expresses themselves. Whether you want to document your life, share knowledge, or start a side business, your phone is your best creative tool — no need to buy a camera or learn Premiere. A single phone is enough to get started. This chapter helps you start from scratch.}

\booksection{I. Think Before You Shoot: Who's Your Audience?}

The first mistake many people make is “picking up the phone and shooting immediately” without knowing who the audience is. Before you start, answer three questions:

\begin{enumerate}
  \item \textbf{Who is my audience?}— Content for kids, peers, or elderly viewers requires completely different styles.
  \item \textbf{What value can I provide?}— Make people laugh (entertainment), learn (knowledge), crave (food), or feel warm (emotion). You need at least one.
  \item \textbf{Which platform can I consistently post on?}— Douyin/Kuaishou (general entertainment), Xiaohongshu (lifestyle), Bilibili (in-depth content), WeChat Channels (social circle). Choose one main platform to focus on — don't distribute across all platforms from the start.
\end{enumerate}

\booksection{II. Phone Camera Basic Settings — 5 Things to Do Before Shooting}

\begin{enumerate}
  \item \textbf{Set Resolution to 1080P, Frame Rate to 30fps}: This is the standard for short video platforms. 4K files are too large and upload slowly. 60fps has no advantage on social platforms and just creates bigger files.
  \item \textbf{Clean the Lens}: The first and most important tip. Wipe the lens with your clothes before shooting — the image immediately becomes noticeably clearer.
  \item \textbf{Lock Focus and Exposure}: Long-press the subject in the viewfinder for 2 seconds until "AE/AF Lock" appears. This prevents brightness fluctuations while moving during the shot.
  \item \textbf{Horizontal or Vertical?}— For Douyin/Kuaishou/WeChat Channels/Xiaohongshu\textbf{shoot vertically}(9:16). For Bilibili/YouTube\textbf{shoot horizontally}(16:9). Decide the platform first, then the orientation.
  \item \textbf{Enable Grid Lines}: Settings > Camera > Grid. The foundation of rule-of-thirds composition, helping you quickly find the visual center.
\end{enumerate}

\booksection{III. Five Camera Movement Techniques — Make Your Footage Come Alive}

The biggest problem for beginner videographers is footage that's too “static” — standing still the entire time. Learn these five basic camera movements, and your footage will instantly feel more professional:

\begin{itemize}
  \item \textbf{Push}: The phone slowly moves toward the subject. Creates tension or intimacy of "approaching."
  \item \textbf{Pull}: The phone slowly moves away from the subject. Pulls from detail to wide shot, establishing the environment.
  \item \textbf{Pan}: The phone stays in place, rotating left or right. Good for landscapes and showing space.
  \item \textbf{Track}: The phone moves parallel, following the subject. Requires steady hands, best with a smartphone gimbal.
  \item \textbf{Follow}: The phone follows behind or in front of the subject. The most commonly used movement in vlogs.
\end{itemize}

\textbf{Stability Is the Top Priority}: Shaky footage looks very unprofessional. Two-handed grip is more stable than one-handed; when walking, bend your knees slightly and use a “ninja walk” to reduce up-and-down bouncing; if you have budget, buy a smartphone gimbal (DJI OM series about 500-800 yuan) — it's a game-changer.

\booksection{IV. Editing Basics — CapCut Is Enough}

JianYing (the domestic version of CapCut) is currently the most beginner-friendly mobile editing tool — free, feature-rich, and template-heavy. Here's the basic workflow:

\begin{enumerate}
  \item \textbf{Import Footage}: Open CapCut > New Project > Select and import your video clips.
  \item \textbf{Trim}: Drag the timeline to trim the wasted footage at the beginning and end of each clip. Principle: keep only the parts with information — shorter is better than longer.
  \item \textbf{Add Transitions}: Tap the "Transition" button between two clips and choose a simple transition effect. Beginners should use "Dissolve" or "Fade to Black" — avoid flashy 3D transitions.
  \item \textbf{Add Subtitles}: CapCut has an "Auto Captions" feature that generates subtitles with one tap. Accuracy is about 90\% — just manually proofread once.
  \item \textbf{Add Music}: CapCut's music library has many royalty-free tracks. Choose a BGM that matches the rhythm and content, and set the volume to about 1/3 of the original video audio.
  \item \textbf{Export}: Export at 1080P, 30fps, bitrate set to "Recommended."
\end{enumerate}

\booksection{V. The Importance of Thumbnails — 70\% of People Decide Whether to Watch Based on the Cover}

A short video's thumbnail is its “storefront” — it determines whether people click. Key points for making thumbnails:

\begin{itemize}
  \item \textbf{Thumbnail Should Have One Sentence}: Write the core content of the video in large text, e.g., "Learn Phone Photo Editing in 3 Minutes." Font should be large, colors bright, with contrast against the background.
  \item \textbf{Thumbnail Must Match Content}: Don't be a "clickbaiter" — if the thumbnail says one thing but the video shows another, viewers will swipe away immediately, and the algorithm will penalize you.
  \item \textbf{Facial Expressions Should Be Vivid}: If the thumbnail has a person, pick a moment with a vivid or exaggerated expression. Human faces attract more clicks than text.
  \item \textbf{Making Thumbnails in CapCut}: In CapCut, select the best frame, add text overlay, and export as an image for the thumbnail.
\end{itemize}

\booksection{VI. Three Content Creation Principles}

\begin{itemize}
  \item \textbf{First 3 Seconds Determine Life or Death}: On short video platforms, viewers decide within 3 seconds whether to swipe or keep watching. The first 3 seconds must have a "hook" — a question, a suspense, or a surprising visual.
  \item \textbf{One Video, One Topic}: The most common mistake beginners make is cramming too much into one video. One video should solve one problem — thoroughness beats comprehensiveness.
  \item \textbf{Persistence Beats Talent}: Algorithmic recommendation requires "warming up the account." The first 20 videos will get almost no views — this is normal. Keep updating 2-3 per week, let the algorithm learn your content direction, and traffic will gradually grow.
\end{itemize}

\begin{tcbpractice}{Hands-On: Shoot a 30-Second Short Video}
\begin{enumerate}
  \item Pick a small topic you know well (e.g., "How to make phone fonts larger").
  \item Shoot 3-5 vertical clips, each 5-10 seconds.
  \item Use CapCut to splice, trim, add subtitles, and add BGM.
  \item Create a thumbnail with large text.
  \item Export and send to a friend first for feedback before publishing.
\end{enumerate}

\end{tcbpractice}

\begin{tcbsummary}{Chapter Summary}
\begin{itemize}
  \item Think before shooting: who's the audience, what value you provide, which platform to post on
  \item 5 things before shooting: 1080P/30fps, clean lens, lock focus/exposure, decide orientation, enable grid
  \item Five camera movements: Push, Pull, Pan, Track, Follow — stability is priority \#1
  \item Edit with CapCut: trim transitions subtitles BGM export — six steps
  \item First 3 seconds are make-or-break, one video one topic, persistence beats talent
\end{itemize}

\end{tcbsummary}
